\section{Verification, provenance, and scope}
\label{sec:provenance}

\subsection{Numerical evaluation of the proved limits}

\Cref{tab:binary} records the binary-field profiles, evaluated from the
proved infinite sums. All entries were recomputed at higher precision and
with larger truncations. They are evaluations of theorems, not estimates
from a simulation.

\begin{table}[htbp]
\centering
\small
\setlength{\tabcolsep}{5pt}
\begin{tabular}{rccccc}
\toprule
$a$ & Full TV & Dimension TV & Bayes recovery & Unique form & Linear uniqueness\\
\midrule
$-2$ & .204664582 & .110014512 & .292186807 & .029601132 & 0\\
$-1$ & .297541183 & .192525397 & .503310958 & .209448088 & 0\\
$0$  & .411496026 & .288788095 & .701564718 & .488389858 & .288788095\\
$1$  & .649349926 & .288788095 & .836030787 & .708770120 & .577576190\\
$2$  & .808320018 & .577576190 & .914018723 & .844641387 & .770101587\\
$3$  & .899761224 & .770101587 & .955969600 & .919767552 & .880116099\\
$4$  & .948740304 & .880116099 & .977719658 & .959230432 & .938790506\\
\bottomrule
\end{tabular}
\caption{Limiting probabilities and total variation distances for $q=2$.
Total variation is a testing advantage, not itself a testing success
probability; equal-prior optimal success is $(1+\mathrm{TV})/2$.}
\label{tab:binary}
\end{table}

For completeness, a convenient evaluation formula at negative as well as
nonnegative offsets is obtained by setting
$r=\lceil\lambda_a\rceil-1$ and writing $B_{m,p}(r)$ for the binomial
cumulative probability $\Pbb(\Bin(m,p)\le r)$. Then
\begin{equation}
 T_q(a)=\sum_{k\ge0}\pi_{a-1}(k)
 \left[B_{m_k,\omega_q}(r)
 -\frac{m_k\omega_q}{\lambda_a}B_{m_k-1,\omega_q}(r-1)\right].
 \label{eq:cdf-evaluation}
\end{equation}
For $m_k=0$, the bracket is one, and the second term is set to zero.
The formula also handles a possible integral threshold because a likelihood
tie contributes zero. It follows from
$T_q(a)=\E(1-Z/\lambda_a)_+$ and the elementary truncated-binomial
first-moment identity.

For $L\ge\max(0,-b)$ with $2L+b+1>0$, the tail of the dimension law obeys
\begin{equation}
 \sum_{k\ge L}\pi_b(k)
 \le\frac{q^{-L(L+b)}}{g_\infty(1-q^{-(2L+b+1)})}.
 \label{eq:tail-compute}
\end{equation}
This bounds the omitted tail in \eqref{eq:cdf-evaluation}; for recovery,
multiply the corresponding null-law bound by $1/\lambda_a$.
For a product $\prod_{j\ge1}(1-x_j)$, its omitted relative factor lies
between $1-\sum_{j>J}x_j$ and one whenever that sum is below one.
For $g_\infty$ the sum is $q^{-J}/(q-1)$, and for the odd product
it is $q^{-(2J+1)}/(1-q^{-2})$. Thus product and summation truncation
errors can be bounded explicitly.

The file \texttt{code/critical\_profiles.py} evaluates 39 profiles for
$q=2,3,5$ and $-4\le a\le8$, checks normalization and the first three
factorial moments, strict information gains, and the $a\ge0$ recovery
identity. It uses 90 decimal digits, then repeats the calculation at 120
digits with longer products and dimension cutoffs. The maximum discrepancy
was below $5\times10^{-88}$. Floating-point roundoff is checked by this
repeat; it is not claimed to be a formal interval-arithmetic certificate.

\subsection{Exact finite-field enumeration}

The standard-library program \texttt{code/exact\_binary\_experiment.py}
computes the actual observation experiment on $\mathbb F_2^4$ using
integers and rational arithmetic. There are 35 two-planes and 28
nondegenerate projective alternating forms. Dynamic programming on row
spans represents every ordered sequence of observations without enumerating
the sequences one by one. It checks all sample lengths $0\le t\le10$,
including the $35^{10}=2{,}758{,}547{,}353{,}515{,}625$ sequences at the
largest length. The largest reachable support contains 2,355 row subspaces.

At $n=4,t=5$ the exact values are
\[
 \E_QZ=\frac{972}{2401},\qquad
 \TV(P,Q)=\frac{199008}{300125},\qquad
 \operatorname{Rec}_{4,5}=\frac{101117}{121500}.
\]
The program verifies the likelihood normalization, exact first moment,
dimension coarsening inequality, recovery ordering, and the exact
existence-test identity. These finite-size numbers differ substantially
from their limits; no unjustified small-$n$ accuracy is being asserted.

A second program, \texttt{code/pencil\_checks.py}, exactly enumerates
independent matrix families for
\[
 (q,n,k)=(2,2,2),(2,4,2),(2,4,3),(3,4,2).
\]
It also records seven reproducible simulations of larger families, with
20,000 samples per case and seed 27112026. Exact enumeration and simulation
are labeled separately. For $q=2,n=4,k=2$, the probability that all three
projective combinations are nonsingular is $336/4096$, whereas the product
of the finite marginal probabilities is $(7/16)^3$. For
$q=2,n=4,k=3$, all seven combinations cannot simultaneously be nonsingular,
although their eventual limiting probability is positive. These checks
guard against accidentally claiming finite-size independence. They do not
replace the proof of \cref{thm:pencil}.

\subsection{Repository snapshot and the precise extension}

The examined repository snapshot is
\begin{center}
\texttt{VladimirReshetnikov/ProveIt}\\
\small\texttt{641cd6296144d7c29bcc6569655b473647428b1b}.
\end{center}
Its GRH project describes a merge of 194 reports, with a separate record
of formalized statements. Its declared input is RH and GRH for Dirichlet
$L$-functions, pinned upstream to
\texttt{72fead102f663d5c7a407964a13dc0cf9d977207}. The vendored development
is described there as unmodified and unbuilt \cite{repo}. We did not rerun
its Lean build or independent kernel checks. The present proof has no
dependency on that input theorem.

The directly relevant source branch is the merged article's section
\emph{Spanning, frames, auxiliary characters and matrix statistics},
especially its questions about an unknown source form. The two immediate
incoming predecessors were verified byte-for-byte against their pinned
Git blobs:
\begin{itemize}
\item \texttt{grh\_hidden\_forms\_and\_filters\_2026-11-27.zip},
which proves the additive sample threshold, exact likelihood, Bayes
identity, and critical chi-square limit \cite{hidden}.
\item \texttt{grh\_finite\_observation\_and\_critical\_recovery.zip},
which proves critical dimension universality, dimension-only testing,
linear recovery, and an observation-overhead law \cite{critical}.
\end{itemize}
The second predecessor explicitly leaves the count of compatible
nondegenerate forms and the full-data optimum open. The new contributions
are full subspace universality, the projective singularity-pattern law,
and the resulting complete critical experiment. The old likelihood and
classical marginal rank law are identified as inputs and rederived where
needed, rather than claimed as new results.

\subsection{Literature and limits of the originality claim}

Classical rectangular and alternating rank laws are treated in
Fulman--Goldstein \cite{fulman}. Exact joint rank counts for
$A,B,A+B$ in alternating ensembles are given by Pott--Schmidt--Zhou
\cite{pott}. Structured quadratic-evaluation rows and their rank
thresholds occur in Cascudo--Cramer--Mirandola--Z\'emor \cite{cascudo};
random rank-one tensors and general-position questions are studied by
Randriambololona \cite{randriam}. The rank-metric density work
\cite{gruica} supplies particularly close context for nonsingularity of
every matrix in a random pencil. The current proof combines a different
row ensemble with a full subspace comparison and a fixed-field,
arbitrary-fixed-dimension alternating pattern theorem.

The targeted search and repository comparison did not locate this complete
critical-law result. That supports treating it as a new result of this
investigation relative to the inspected material; it does not establish
worldwide publication priority. The proof has been independently checked
within this task, but has not been externally peer reviewed or formalized
in Lean.

The asymptotic parameters are explicit: $q$, the critical offset $a$, and
the dimension in the standalone pencil theorem are fixed. No uniform
claim is made when these grow with $n$. \Cref{ker:finite-TV} supplies a
computable finite comparison for the kernel step, but the combined
optimal-risk theorem is stated as a limit rather than an optimized
finite-size error estimate. These boundaries do not affect any of the
displayed exact finite identities.
